\documentclass[a4paper]{article}

\usepackage{graphicx}
\usepackage{amsmath,amssymb}
\usepackage{minted}
\usepackage{multirow}
\usepackage{float}
\usepackage{algorithm}
\usepackage{algpseudocode}
\usepackage{todonotes}
\usepackage{forest}
\usepackage{xstring}
\usepackage{xcolor}
\usepackage[most]{tcolorbox}
\usepackage{xparse}
\usepackage{natbib}

\title{Navigator refresh after rename and move}
\date{2026-09-07}
\author{Mohammad Rahmani}

\begin{document}
\maketitle

\section{Problem}
Navigator could keep stale rows after a file or folder had already been renamed
or moved successfully on disk. The filesystem operation completion path
requested a proxy refresh, but the proxy model intentionally ignored
\texttt{invalidateFilter()}, so the old proxy mapping could remain visible.

A second ownership problem allowed a later workspace interaction layer to
disconnect and reconnect the Navigator context-menu signal even though the
multi-path context controller is the canonical menu owner.

\section{Fix}
\begin{itemize}
   \item Explicit filesystem mutation refreshes now rebuild the
   \texttt{QSortFilterProxyModel} mapping from the current
   \texttt{QFileSystemModel} state.
   \item Existing create, rename, move, delete, and paste completion paths keep
   using the same refresh entry point, so no parallel refresh mechanism is
   introduced.
   \item The workspace interaction controller no longer owns or reconnects the
   Navigator context menu.
   \item \texttt{NavigatorMultiPathContextController} remains the single
   right-click menu owner, including
   \textbf{Copy / Share $\rightarrow$ Copy from Default Path}.
\end{itemize}

\section{Regression coverage}
Regression checks verify real proxy invalidation after filesystem mutations,
existing move/delete refresh calls, single context-menu ownership, and the
presence of \textbf{Copy from Default Path}.

\end{document}
